\chapter{Допуск котангенциального удвоенного колчанного родителя}
\label{ch:v8-horizontal-phase-cotangent-doubled-quiver-parent-admission-gate}

Прежде чем вводить новый родитель, необходимо сопоставить его с уже
пройденными уровнями. В Томе III BF/AKSZ-форма фиксировала скобку, но не
положительную метрику и не поляризацию. В Томе V предпроективное отношение
точно воспроизводило разность произведений стрелок, но положительная
инволюция и физическая группа не выводили требуемое отображение момента.
Предыдущий гейт впервые дал невырожденную комплексную поляризацию полного
сбалансированного носителя. Поэтому старый вопрос можно повторить на новом
уровне, не объявляя прежний no-go отменённым.

Пусть $(\mathcal T_{26},\Omega_0)$ --- построенное симплектическое
представление, а $\rho(X_a)$ --- четырнадцать образующих
комплексифицированного калибровочного действия. Из инвариантности формы
следует
\begin{equation}
 K_a:=\Omega_0\rho(X_a),qquad K_a^{\mathsf T}=K_a.
 \label{eq:v8-cotangent-moment-quadratic-matrices}
\end{equation}
Тем самым каноническое квадратичное отображение момента определено формулой
\begin{equation}
 \mu_a(x)=\frac12x^{\mathsf T}K_ax
 =\frac12\Omega_0\!\left(x,\rho(X_a)x\right).
 \label{eq:v8-cotangent-moment-map}
\end{equation}
В отличие от раннего вещественного $SO(3)$-среза оно не равно нулю.
Например, на одном слабом котангенциальном базисном паре точный вектор
компонент равен
\begin{equation}
 \mu(e_1+e_9)=(-1,0,\ldots,0,-1/2)^{\mathsf T}.
 \label{eq:v8-cotangent-nonzero-moment-witness}
\end{equation}

Однако четырнадцать записанных компонент имеют одну точную центральную
зависимость. Ранг семейства квадратичных матриц равен
\begin{equation}
 \rank\operatorname{span}\{K_a\}_{a=1}^{14}=13.
 \label{eq:v8-cotangent-moment-span-rank}
\end{equation}
Коэффициенты зависимости в порядке четырёх матричных единиц слабого блока,
девяти цветного блока и гиперзаряда равны
\begin{equation}
 \left(-\frac12,0,0,-\frac12,
 -\frac23,0,0,0,-\frac23,0,0,0,-\frac23,1\right).
 \label{eq:v8-cotangent-central-generator-relation}
\end{equation}
Это не зануляет отображение момента, а удаляет одну избыточную центральную
координату.

Теперь рассмотрим каноническое действие фазы слоя котангенциальных
волокон. На тринадцати координатах исходных стрелок оно действует как $z$,
а на тринадцати двойственных координатах как $z^{-1}$:
\begin{equation}
 S_z=\operatorname{diag}(zI_{13},z^{-1}I_{13})
 \quad\text{в попарно упорядоченном базисе},qquad z\in\mathbb C^\times.
 \label{eq:v8-cotangent-fibre-phase}
\end{equation}
В используемом изотипическом порядке та же диагональ переставлена вместе с
базисом. Точная проверка даёт
\begin{equation}
 S_z^{\mathsf T}\Omega_0S_z=\Omega_0,qquad
 [S_z,\rho(X_a)]=0,qquad
 S_z^{\mathsf T}K_aS_z=K_a.
 \label{eq:v8-cotangent-phase-invariance}
\end{equation}
Следовательно,
\begin{equation}
 \mu(S_zx)=\mu(x),qquad
 V_{\lambda,\zeta}(S_zx)=V_{\lambda,\zeta}(x),qquad
 V_{\lambda,\zeta}(x)=\lambda\|\mu(x)-\zeta\|^2.
 \label{eq:v8-cotangent-parent-phase-blindness}
\end{equation}
Любая функция только от отображения момента остаётся слепой к этой фазе.

Кроме того, симплектическая форма не выбирает положительную норму на
двойственной алгебре, общий коэффициент $\lambda$ или уровень стабильности
$\zeta$. Уже значение в начале координат равно
$V_{\lambda,\zeta}(0)=\lambda\|\zeta\|^2$, поэтому разные $\zeta$ задают
разные родительские теории, а не калибровочные записи одного объекта.

Это согласуется со стандартной картиной: удвоение колчана реализует
котангенциальный симплектический носитель, а колчанные многообразия строятся
редукцией по уровню отображения момента
\cite{v8-cotangent-crawley-boevey,v8-cotangent-harada-wilkin}.
Литература предоставляет форму конструкции, но не выбирает проектные
$\lambda$, $\zeta$ и отождествление новой фазы с наблюдаемой модой.

\begin{theorem}[Ненулевое отображение момента полного носителя]
На $\mathcal T_{26}$ каноническая симплектическая формула задаёт ненулевое
тринадцатимерное отображение момента физического калибровочного действия.
Тем самым ранний нулевой результат на вещественном $SO(3)$-срезе не
переносится на новый комплексный носитель.
\end{theorem}

\begin{nogocriterion}[No-go канонического котангенциального подъёма фазы]
Каноническая фаза котангенциального волокна сохраняет $\Omega_0$, коммутирует
с калибровочным действием и оставляет каждую компоненту $\mu$ неизменной.
Поэтому ни $\|\mu\|^2$, ни сдвинутая норма $\|\mu-\zeta\|^2$ не создают
массу горизонтальной фазе. Нарушить эту симметрию можно лишь дополнительной
комплексной структурой, метрикой или взаимодействием, не являющимся функцией
одного отображения момента.
\end{nogocriterion}

\begin{protocol}
\begin{enumerate}
 \item симплектический носитель ранга $26$: \textbf{имеется};
 \item квадратичные матрицы $K_a$: \textbf{симметричны};
 \item размерность отображения момента: \textbf{$13$};
 \item ненулевой точный свидетель: \textbf{имеется};
 \item центральных зависимостей образующих: \textbf{$1$};
 \item каноническая котангенциальная фаза симплектична: \textbf{да};
 \item все компоненты $\mu$ видят эту фазу: \textbf{нет};
 \item норма отображения момента поднимает фазу: \textbf{нет};
 \item положительная метрика на алгебре выведена из $\Omega_0$: \textbf{нет};
 \item коэффициент $\lambda$ и уровень $\zeta$: \textbf{свободны};
 \item новый колчанный родитель канонически выбран: \textbf{нет};
 \item следующий гейт: \textbf{селектор совместимой комплексной структуры
 и положительной метрики}.
\end{enumerate}
\end{protocol}

Машинный сертификат сохранён в
\path{s2t_v8_horizontal_phase_cotangent_doubled_quiver_parent_admission_gate_results.json}.

\begin{thebibliography}{9}
\bibitem{v8-cotangent-crawley-boevey}
W.~Crawley-Boevey,
\emph{Poisson structures on moduli spaces of representations},
J. Algebra 325 (2011), 205--215; arXiv:math/0502301.
\bibitem{v8-cotangent-harada-wilkin}
M.~Harada, G.~Wilkin,
\emph{Morse theory of the moment map for representations of quivers},
Geom. Dedicata 150 (2011), 307--353; arXiv:0807.4734.
\end{thebibliography}